\section{Second order at every rank}
\label{sec:second}

This section is the algebraic base of the whole program. Everything in
it is proof-checked, and its conclusion --- that the band exists for
$N \ge 3$ and not for $N = 2$, with an exact amplitude --- follows from
representation theory with no input from any expansion.

\subsection{The shared-link weights are Weingarten values}

Two plaquettes contributing at second order either share a link or do
not. The amplitude for the shared-link case was for a long time taken
in this program with an isotropy premise: that the four contributing
representation channels carry weights in a fixed ratio. That premise
is not needed. It is a consequence of the order-two Weingarten values.

\begin{proposition}[The shared-link weights, from representation theory alone]
\label{prop:weingarten}
Let two plaquettes share exactly one link. Then:
\begin{enumerate}
  \item The six non-shared links collapse. Each plaquette contributes a
    product of three independent Haar links, and a product of
    independent Haar matrices is Haar. So the two-plaquette amplitude
    reduces to $\Tr(A U)\,\Tr(B U^{\pm 1})$ with $A$ and $B$ independent
    Haar, and integrating them out leaves a pure degree-$(2,2)$ moment of
    the shared link $U$.
  \item Two such moments settle both families. With
    $M_{\mathrm{direct}} = N^2$ and $M_{\mathrm{cross}} = N$, the like
    family splits as $(N+1)/(2N)$ and $(N-1)/(2N)$; and the singlet
    component of $U_{ij}\overline{U_{lk}}$ is
    $\delta_{il}\delta_{jk}/N$, of squared norm $1$, giving $1/N^2$.
\end{enumerate}
All four weights are exactly $d_R/N^2$, where $d_R$ is the dimension of
the channel.
\end{proposition}
\gidc{the shared-link weights are Weingarten, not an isotropy assumption}{workhouse verify --only 'the shared-link weights are Weingarten, not an isotropy assumption'}

\noindent\Tone{}. The relevant Weingarten pair is the inverse of the
$S_2$ Gram matrix and the index sums are explicit, so the check imports
nothing from the corpus. Registered as the discharge of the gap that
had carried this premise as the one unproved input of the second-order
chain; it was opened and closed on the same day, 28~August 2026, once
the question was asked in the right form.

\subsection{The four channels in closed form}

Write $C_F = (N^2-1)/(2N)$ for the fundamental Casimir. Each channel
$R$ contributes a resolvent weight $-(d_R/N^2)/(C_F + \tfrac12 \Delta C_2(R))$,
and all four evaluate in closed form.

\begin{theorem}[Per-channel resolvent weights]
\label{thm:channels}
\begin{align}
  w_{\mathbf 1}   &= -\frac{1/N^2}{C_F}
    &&= -\frac{2}{N(N^2-1)},
  \\[3pt]
  w_{\mathrm{Adj}} &= -\frac{(N^2-1)/N^2}{C_F + N/2}
    &&= -\frac{2(N^2-1)}{N(2N^2-1)},
  \\[3pt]
  w_{\mathrm{Sym}} &= -\frac{N(N+1)/2N^2}
                        {C_F + \frac{(N-1)(N+2)}{2N}}
    &&= -\frac{N+1}{(N-1)(2N+3)},
  \\[3pt]
  w_{\Lambda} &= -\frac{N(N-1)/2N^2}
                    {C_F + \frac{(N+1)(N-2)}{2N}}
    &&= -\frac{N-1}{(N+1)(2N-3)} .
\end{align}
\end{theorem}
\gid{LEAN:resolventWeight\_singlet, LEAN:resolventWeight\_adjoint, LEAN:resolventWeight\_sym, LEAN:resolventWeight\_antisym}

\noindent\Tzero{} --- all four, at
\srcpath{lean/Workhouse/Basic.lean:72,80,89,104}. Reproduce with
\texttt{make lean}.

\subsection{The two family sums}

The channels group into a \emph{mixed} family $\{\mathbf 1, \mathrm{Adj}\}$
and a \emph{like} family $\{\Lambda^2 F, \mathrm{Sym}^2 F\}$, and these
are the antiparallel and parallel channel sums.

\begin{theorem}[Family sums]
\label{thm:families}
\begin{equation}
  A_N \;=\; w_{\mathbf 1} + w_{\mathrm{Adj}}
    \;=\; -\,\frac{2N^3}{(N^2-1)(2N^2-1)},
  \qquad
  B_N \;=\; w_{\mathrm{Sym}} + w_{\Lambda}
    \;=\; -\,\frac{4N(N^2-2)}{(N^2-1)(4N^2-9)} .
\end{equation}
\end{theorem}
\gid{LEAN:mixed\_family\_sum, LEAN:like\_family\_sum}

\noindent\Tzero{}, at \srcpath{lean/Workhouse/Basic.lean:118,129}. The
mixed sum is obtained from the dimension and Casimir table rather than
by transcription from any source document, which is the point: the
family sums are re-derived, not quoted.

\subsection{The hopping amplitude}

\begin{theorem}[The all-rank second-order hopping amplitude]
\label{thm:tN}
\begin{equation}
  t_N \;=\; B_N - A_N
  \;=\; \frac{2N\,(N^2-4)}{(N^2-1)(2N^2-1)(4N^2-9)} .
  \label{eq:tN}
\end{equation}
It is strictly positive for every $N \ge 3$, and
\begin{equation}
  t_2 = 0, \qquad t_3 = \tfrac{5}{612}.
\end{equation}
\end{theorem}
\gidc{CONST:t\_N}{workhouse why t\_N}

\noindent\Tzero{} for the rank law and for both special values
(\srcpath{lean/Workhouse/Basic.lean}: \texttt{rank\_law\_numerator},
\texttt{hopping\_two}, \texttt{hopping\_three}). Clearing denominators,
the channel difference has numerator exactly $2N(N^2-4)$.

Two consequences deserve to be stated separately, because both are
structural rather than numerical.

\begin{corollary}[$\SU(2)$ has no band at this order]
\label{cor:su2}
The factor $N^2 - 4$ in \eqref{eq:tN} vanishes at $N = 2$. The exclusion
of $\SU(2)$ is therefore not a smallness but an identity: the parallel
and antiparallel channel sums coincide exactly at rank two.
\end{corollary}
\gid{LEAN:hopping\_two}

\begin{corollary}[The planar approach is from below]
\label{cor:deficit}
With $D = (N^2-1)(2N^2-1)(4N^2-9)$,
\begin{equation}
  D - 4N^3 \cdot 2N(N^2-4) \;=\; 2N^4 + 31N^2 - 9 \;>\; 0
  \qquad (N \ge 1),
\end{equation}
so $4N^3 t_N = 1 - (2N^4+31N^2-9)/D < 1$. Hence
$t_N < 1/(4N^3)$ for every $N$, and $t_N \sim 1/(4N^3)$: the
$N^{-3}$ cancellation is exact and what it leaves is a strictly positive
deficit.
\end{corollary}
\gid{LEAN:hopping\_deficit\_numerator}

\noindent\Tzero{}, at \srcpath{lean/Workhouse/Basic.lean:48}.
Corollary~\ref{cor:deficit} is the first appearance of the rank grading
that \S\ref{sec:grading} develops: the leading power of $N$ is fixed by
a cancellation, and the surviving term is a polynomial one can write
down.

\subsection{Independent corroboration of the channel weights}

The weights of Proposition~\ref{prop:weingarten} have an external
provenance, from work written without knowledge of this program.

An efficient-truncation study of $\SU(N_c)$ lattice gauge theory for
quantum simulation~\cite{CBB_2026} gives, in its Appendix~D, the exact
finite-rank plaquette matrix element
\begin{equation}
  \abs{M_\rho} = \sqrt{\frac{\dim\rho}{\dim A \,\dim R}} .
\end{equation}
Taking both factors fundamental gives $\abs{M_\rho}^2 = d_\rho/N^2$,
which is \emph{identically} the numerator of the weight formula used
here. The four shared-link channel weights, and with them $A_N$, $B_N$
and $t_N = B_N - A_N$, therefore rest on a matrix element an
independent group wrote down for an unrelated purpose. That truncation
is also the first published one retaining all four shared-link
irreducible representations.

The agreement is evidence rather than bookkeeping precisely because
neither side knew of the other. It does not make $t_N$ published: the
rank law \eqref{eq:tN}, the vanishing at $N = 2$, and the deficit of
Corollary~\ref{cor:deficit} are assembled here.

A caution attaches in the other direction. The leading-large-$N_c$
qutrit plaquette Hamiltonian~\cite{CB_2024} explicitly decouples the
charge-odd combination, so the sector in which $t_N$ lives is inert in
that model --- effectively $t_3 = 0$ at leading order in large $N_c$,
against the finite-$N$ channel-complete value $5/612$. That paper is
recorded in this program's index as not yet obtained, and the statement
here is taken from imported bridge documents rather than checked
against the paper itself.

\subsection{What this section does not establish}

\notasserted{Convergence of the strong-coupling series; anything about
the spectrum at finite $u$ beyond second order; and nothing about
$\SU(2)$ at higher orders, where Corollary~\ref{cor:su2} does not
apply. Theorem~\ref{thm:tN} is a statement about one coefficient of a
formal expansion, established exactly.}
